
\documentclass[11pt]{article}
\usepackage{amsmath,amssymb,amsthm,mathtools}
\usepackage{enumitem}
\usepackage{booktabs}
\usepackage[margin=1in]{geometry}
\usepackage{hyperref}

\newtheorem{theorem}{Theorem}
\newtheorem{lemma}{Lemma}
\newtheorem{definition}{Definition}
\newtheorem{proposition}{Proposition}
\newtheorem{corollary}{Corollary}
\newtheorem{remark}{Remark}

\title{Step 146: Trace-Class Residual Ledger Theorem for $K_R$}
\author{RATCHET / Formed-Layer Membrane RH Thread}
\date{}

\begin{document}
\maketitle

\section{Purpose}
Step 145 instantiated the completed residual carrier
\[
H_R=\overline{\operatorname{span}\{\operatorname{Ran}\Pi_{Y_a}\mathfrak B_{\ell,a}:\ell\in\mathcal L\setminus\{0\}\}}
\subseteq Y_a\subset L_a,
\]
where $Y_a$ is Burnol's closed span of zero-evaluator vectors and $P_a=Y_a^\perp$ is the co-Poisson complement.
Step 146 asks whether the completed residual ledger
\[
K_R=\Pi_RK_Z\Pi_R
\]
is trace-class on $H_R$.  The answer is: not automatically.  A trace-class residual ledger is available only after an explicit weight schedule and evaluator-norm summability record are declared.

\section{Zero-evaluator readout}
Let
\[
\mathcal Z=\{(\rho,k):\zeta(\rho)=0,\;0\le k<m_\rho\}
\]
be the nontrivial zero/multiplicity index set.  Burnol supplies evaluator vectors
\[
Y^a_{\rho,k}\in L_a,
\qquad
\langle f,Y^a_{\rho,k}\rangle=M(f)^{(k)}(\rho),
\]
for $\rho\neq 0,1$.  Let $\Pi_R:L_a\to H_R$ denote the orthogonal projection.  Define the off-critical readout
\[
q(\rho):=\Re\rho-\frac12.
\]

\begin{definition}[Weighted residual analysis operator]
Let $\omega_{\rho,k}>0$ be an upstream-declared weight schedule.  Define
\[
D_R^{\omega}:H_R\to \ell^2(\mathcal Z),
\qquad
(D_R^{\omega}h)_{\rho,k}
=\sqrt{\omega_{\rho,k}}\,q(\rho)\,\langle h,\Pi_RY^a_{\rho,k}\rangle.
\]
The weighted residual ledger is
\[
K_R^{\omega}:=(D_R^{\omega})^*D_R^{\omega}.
\]
Equivalently,
\[
K_R^{\omega}
=\sum_{(\rho,k)\in\mathcal Z}
\omega_{\rho,k}q(\rho)^2
\,|\Pi_RY^a_{\rho,k}\rangle\langle\Pi_RY^a_{\rho,k}|,
\]
where the sum is understood in the monotone positive-form sense.
\end{definition}

\section{Trace-class criterion}
\begin{theorem}[Trace-class residual ledger criterion]
The following are equivalent for the positive ledger defined above:
\[
D_R^{\omega}\in\mathcal S_2(H_R,\ell^2(\mathcal Z)),
\]
\[
K_R^{\omega}\in\mathcal S_1(H_R),
\]
and
\[
\boxed{
\sum_{(\rho,k)\in\mathcal Z}
\omega_{\rho,k}q(\rho)^2
\|\Pi_RY^a_{\rho,k}\|_{H_R}^2<\infty.
}
\]
When these conditions hold,
\[
\boxed{
\operatorname{tr}K_R^{\omega}
=\sum_{(\rho,k)\in\mathcal Z}
\omega_{\rho,k}q(\rho)^2
\|\Pi_RY^a_{\rho,k}\|_{H_R}^2.
}
\]
\end{theorem}

\begin{proof}
For a rank-one positive operator $|v\rangle\langle v|$, the trace is $\|v\|^2$.  The monotone partial sums
\[
K_F^\omega=\sum_{(\rho,k)\in F}\omega_{\rho,k}q(\rho)^2|\Pi_RY^a_{\rho,k}\rangle\langle\Pi_RY^a_{\rho,k}|
\]
are positive and have traces equal to the corresponding finite sums.  The positive series defines a trace-class operator if and only if these traces are uniformly bounded; the trace is then the monotone limit.  Equivalently, the analysis operator has Hilbert--Schmidt norm equal to the displayed sum.
\end{proof}

\section{Tail promotion}
Let $P_N\uparrow I_{H_R}$ be the finite residual-window ladder from Step 145: shift modes, zero-height cutoff, and source-readable $\Omega$-compatible dictionary windows.  Set
\[
T_N:=I_{H_R}-P_N.
\]

\begin{theorem}[Trace-class tail promotion]
If $K_R^{\omega}\in\mathcal S_1(H_R)$ and $P_N\uparrow I_{H_R}$ strongly, then
\[
\boxed{
\operatorname{tr}(T_NK_R^{\omega}T_N)\to0.
}
\]
In particular, the completed residual-tail term in Step 144 is paid.
\end{theorem}

\begin{proof}
For a positive trace-class operator $K$, strong convergence $P_N\uparrow I$ implies $(I-P_N)K^{1/2}\to0$ in Hilbert--Schmidt norm.  Hence
\[
\operatorname{tr}((I-P_N)K(I-P_N))
=\|(I-P_N)K^{1/2}\|_{\mathcal S_2}^2\to0.
\]
\end{proof}

\section{Separation and nonclaim}
Trace-class weighting does not destroy separation provided every weight is strictly positive.

\begin{proposition}[Weighted separation]
Assume $\omega_{\rho,k}>0$ for all $(\rho,k)$ and $\|\Pi_RY^a_{\rho,k}\|>0$ for every residual zero-evaluator direction under consideration.  Then
\[
\operatorname{tr}K_R^{\omega}=0
\]
implies
\[
q(\rho)=0
\]
for every zero direction visible in $H_R$.  Thus the weighted ledger separates off-critical residual zero directions.
\end{proposition}

\begin{remark}[Scope]
This is not yet RH.  It is the residual-carrier statement:
\[
\text{no visible residual off-critical zero direction in }H_R.
\]
To upgrade it to the zeta ledger one still needs the already-named adequacy/separation records: every off-critical zero probe must be visible in the residual carrier or paid by a declared $\Xi$-residual.
\end{remark}

\section{Unweighted ledger warning}
The unweighted ledger
\[
\sum_{(\rho,k)}q(\rho)^2|\Pi_RY^a_{\rho,k}\rangle\langle\Pi_RY^a_{\rho,k}|
\]
is not trace-class certified.  Since the number of zeros up to height $T$ grows on the order of $T\log T$, and the evaluator norms need not decay, no unweighted trace-class claim is admissible without a separate zero-evaluator norm theorem.

\section{Completed squeeze}
Combining Step 144 with the trace-class ledger gives the completed residual squeeze
\[
\operatorname{tr}(G_RK_R^\omega)
\le
\frac{C_{\rm src}}{\Lambda_N^\Omega}
+
\operatorname{tr}((I-P_N)K_R^\omega(I-P_N)).
\]
Therefore the residual zero mass vanishes if
\[
\Lambda_N^\Omega\to\infty
\]
and the trace-class tail theorem applies.

\section{Next obligation}
The active analytic obligation is now a lawful weight/evaluator-norm envelope:
\[
\sum_{(\rho,k)}\omega_{\rho,k}q(\rho)^2\|\Pi_RY^a_{\rho,k}\|^2<\infty,
\qquad
\omega_{\rho,k}>0.
\]
The next step should choose an upstream-declared schedule $\omega_{\rho,k}$ using Burnol evaluator-norm growth and zero-counting, then verify that it remains separating and source-compatible.

\end{document}
